(1) 约化质量 $\mu=m_1m_2/(m_1+m_2)$，质心系相对拉氏量 $L=\mu(\dot r^2+r^2\dot\phi^2)/2-k(r-r_0)^2/2$。角坐标 $\phi$ 为循环坐标，故 $\ell=\partial L/\partial\dot\phi=\mu r^2\dot\phi$ 为常数；质心的匀速运动可与相对动力学完全分离。
(2) 消去 $\dot\phi$ 得 $U_{\rm eff}(r)=k(r-r_0)^2/2+\ell^2/(2\mu r^2)$，圆轨道条件 $U_{\rm eff}'(r_c)=k(r_c-r_0)-\ell^2/(\mu r_c^3)=0$。左侧函数导数 $k+3\ell^2/(\mu r^4)>0$，在 $r\to0^+$ 为负无界、$r\to\infty$ 为正无界，故唯一正根且必满足 $r_c>r_0$。
(3) 二阶展开有效势得 $\omega_r^2=U_{\rm eff}''(r_c)/\mu=k/\mu+3\ell^2/(\mu^2r_c^4)$；圆周角速度 $\Omega=\ell/(\mu r_c^2)$。因此 $\omega_r^2=k/\mu+3\Omega^2$，而利用圆轨道条件还可写 $\Omega^2=k(r_c-r_0)/(\mu r_c)$；比例 $\omega_r/\Omega=\sqrt{3+(k\mu r_c^4)/(\ell^2)}$，确为无量纲。
(4) $\ell\to0$ 时屏障消失、$r_c\to r_0$、$\Omega\to0$，径向频率趋 $\sqrt{k/\mu}$；大角动量时 $r_c\gg r_0$，由 $kr_c^4\simeq\ell^2/\mu$ 得 $\Omega^2\to k/\mu$、$\omega_r^2\to4k/\mu$。总相对能量 $E=\mu\dot r^2/2+U_{\rm eff}(r)$，允许区必须满足 $E\ge U_{\rm eff}$，两侧转向点由等号确定。
